\documentclass[conference]{IEEEtran}
\IEEEoverridecommandlockouts
% The preceding line is only needed to identify funding in the first footnote. If that is unneeded, please comment it out.
\usepackage{cite}
\usepackage{amsmath,amssymb,amsfonts}
\usepackage{algorithmic}
\usepackage{graphicx}
\usepackage{textcomp}
\usepackage{xcolor}
\usepackage{booktabs}
\usepackage{multirow}
\usepackage{url}

\def\BibTeX{{\rm B\kern-.05em{\sc i\kern-.025em b}\kern-.08em
    T\kern-.1667em\lower.7ex\hbox{E}\kern-.125emX}}

\begin{document}

\title{VeriByte: Design, Formalization, and Empirical Soundness Evaluation of a Secure Bytecode Virtual Machine with Static Dataflow Verification}

\author{\IEEEauthorblockN{Dhyanam M. N.}
\IEEEauthorblockA{\textit{Department of Computer Science and Engineering} \\
\textit{School of Engineering and Technology}\\
Email: student@university.edu}
\and
\IEEEauthorblockN{Faculty Advisor / Guide}
\IEEEauthorblockA{\textit{Department of Computer Science and Engineering} \\
\textit{School of Engineering and Technology}\\
Email: faculty@university.edu}
}

\maketitle

\begin{abstract}
Virtual machines executing intermediate bytecode representations---exemplified by the Java Virtual Machine (JVM), WebAssembly (Wasm) runtimes, and Linux extended Berkeley Packet Filters (eBPF)---form foundational infrastructure for sandboxed and portable code execution. However, when bytecode originates from untrusted compilers, malicious adversaries, or corrupted transport channels, direct execution induces fatal vulnerabilities including stack underflow/overflow, type confusion, out-of-bounds local access, and control-flow hijacking. Traditional sandboxing mitigates these threats through continuous runtime boundary checks, introducing substantial instruction dispatch overhead (often 200\%--500\%). 

This paper introduces \textbf{VeriByte}, an end-to-end secure programming language toolchain and stack-based virtual machine governed by an ahead-of-time static bytecode verifier based on abstract interpretation. The verifier models program execution over an abstract domain of typed symbolic operand stacks and definitely assigned variable frames across a monotone dataflow framework. We formalize thirteen comprehensive safety rules encompassing control-flow validity, operand-type safety, frame convergence, definite assignment, and call/return isolation. We demonstrate an attack-and-defence paradigm across 26 handcrafted exploit binaries and evaluate empirical soundness via differential mutation fuzzing across 15,000 structural mutants and 10,000 raw-byte corrupted binaries under an active runtime checking oracle. Over 25,000 stress inputs, the verifier exhibited zero false rejections on valid programs and zero soundness violations across all accepted binaries, while preventing 9,421 catastrophic unhandled crashes. Load-time verification imposes strictly linear $\mathcal{O}(N)$ scaling with an overhead below 0.5\% of execution time.
\end{abstract}

\begin{IEEEkeywords}
Bytecode Verification, Abstract Interpretation, Type Safety, Virtual Machine, Static Dataflow Analysis, Soundness Oracle, Differential Fuzzing.
\end{IEEEkeywords}

\section{Introduction}
Modern software deployment relies heavily on platform-independent bytecode architectures. Compiling source code into an intermediate bytecode distributed to a client-side or in-kernel virtual machine (VM) decouples language implementations from hardware targets while theoretically establishing a sandboxed isolation boundary. Representative architectures span the Java Virtual Machine \cite{lindholm2014java}, the Common Language Runtime (CLR) \cite{meijer2002technical}, browser-embedded WebAssembly engines \cite{haas2017bringing}, and kernel-space eBPF execution engines \cite{ebpf2023linux}.

In standard compiler curricula and baseline runtime engines, compiler design ends at syntax-directed code generation, leaving the execution engine to implicitly trust the generated bytecode. In adversarial environments, this implicit trust model fails catastrophically. A binary executable can be hand-crafted in binary hex editors, emitted by compromised third-party compilers, or corrupted during network transit. When an untrusted bytecode binary containing invalid control-flow targets or mismatched stack operands enters a native VM loop, it triggers:
\begin{enumerate}
    \item \textbf{Stack Underflow and Overflow:} Popping from an empty evaluation stack reads adjacent host memory; pushing beyond capacity corrupts activation frames.
    \item \textbf{Control-Flow Hijacking:} Jump targets pointing into the middle of multi-byte instruction encodings cause data operands to be decoded as opcodes, leading to arbitrary code execution.
    \item \textbf{Type Confusion:} Arithmetic instructions executing against boolean flags or object references bypass high-level type safety invariants.
    \item \textbf{Uninitialized Local Access:} Reading unassigned local variable slots leaks unscrubbed host memory registers.
\end{enumerate}

To counter these vulnerabilities without paying the steep runtime performance penalty of defensive interpreter traps (which require checking stack pointers and dynamic types at every fetch-decode-execute step), modern architectures enforce static bytecode verification prior to program instantiation. 

In this work, we present \textbf{VeriByte}, a complete, fully verifiable stack-based compilation and execution pipeline designed as a trust boundary. VeriByte implements an imperative source language (\textit{MiniLang}), a two-pass semantic type analyzer, a 27-opcode stack-based binary bytecode generator, a bytecode optimizer with constant folding and jump threading, a static bytecode verifier employing Kildall's worklist abstract interpretation algorithm \cite{kildall1973unified}, and a high-performance virtual machine interpreter equipped with a checked-mode verification oracle.

The primary contributions of this paper are:
\begin{itemize}
    \item \textbf{System Architecture as a Security Boundary:} We establish an end-to-end verifiable pipeline enforcing the principle that no untrusted bytecode is ever admitted to the runtime interpreter without preceding inductive safety verification.
    \item \textbf{Thirteen Formalized Safety Rules:} We formalize and implement thirteen static safety invariants spanning bytecode structure, control-flow graphs (CFGs), operand-type lattices, join-point stack unification, and call/return typing.
    \item \textbf{Attack-and-Defence Corpus:} We validate verifier precision against a handcrafted corpus of 26 exploit binaries designed to isolate each rule violation, juxtaposing unverified execution crashes against clean verifier rejections.
    \item \textbf{Checked Soundness Oracle and Differential Fuzzing:} We evaluate verifier soundness via a 15,000-mutant structural fuzzing campaign and a 10,000-mutant raw-byte serializer fuzzing campaign. We validate the fuzzer's sensitivity through an intentional verifier weakening meta-test.
    \item \textbf{Empirical Performance Profiling:} We demonstrate that verification scales linearly ($\approx 180\text{ KB/s}$) with an execution overhead under 0.5\%, confirming the practical viability of load-time verification.
\end{itemize}

\section{Literature Survey}
The theoretical foundations of static bytecode verification originate in abstract interpretation, type systems, and monotone dataflow frameworks.

\subsection{JVM Bytecode Verification and Stack Maps}
Early iterations of the Java Virtual Machine Specification \cite{lindholm2014java} implemented dataflow-based type inference to reconstruct operand stack shapes at runtime. In seminal work, Leroy \cite{leroy2003java} formalized Java bytecode verification using abstract interpretation over an abstract machine, establishing proof obligations for type safety and progress-preservation properties. However, because multi-path type inference exhibits polynomial computational complexity, Sun Microsystems introduced the \texttt{StackMapTable} attribute in Java ME and Java 6+. Rather than inferring types from scratch, the modern JVM validator verifies pre-computed type annotations in linear time. VeriByte deliberately adopts the more academically rich full-inference abstract interpretation approach, solving the fixpoint over loop back-edges without requiring external compiler-assisted stack maps.

\subsection{WebAssembly Validation}
WebAssembly (Wasm) was designed specifically to provide memory-safe, near-native execution inside web browsers \cite{haas2017bringing}. Unlike arbitrary stack machines, Wasm enforces structured control flow (avoiding arbitrary goto/jump instructions). This enables single-pass linear validation where operand stacks are validated syntactically using typed block signatures. While structured control flow simplifies validation, legacy and low-level architectures rely on unstructured jump primitives. VeriByte tackles the more generalized, unstructured control-flow domain where jumps can target arbitrary byte offsets, necessitating explicit CFG decoding and merge reconciliation.

\subsection{Linux Kernel eBPF Verification}
Extended Berkeley Packet Filters (eBPF) permit user-space applications to execute sandboxed code inside the Linux kernel \cite{ebpf2023linux, gershuni2019simple}. Because kernel panics are unacceptable, the kernel's in-tree verifier performs symbolic execution to simulate all possible execution paths, validating that registers contain safe data types, pointer arithmetic stays within allocated buffer boundaries, and programs terminate within bounded instruction limits. VeriByte adopts a comparable threat model: untrusted binaries submitted from external sources must pass static analysis before admission to the runtime.

\subsection{Abstract Interpretation Foundations}
Cousot and Cousot \cite{cousot1977abstract} established abstract interpretation as a unified lattice model for approximating the semantics of programming languages. In VeriByte, the concrete execution domain (holding 32-bit integer values and boolean bits) is abstracted into a finite type lattice ($\mathbb{L} = \{\bot, \text{int}, \text{bool}, \top\}$), augmented with definite assignment tracking states ($\{\text{UNSET}, \text{CONFLICT}\}$). Monotonicity over this bounded-height lattice guarantees convergence to a fixed point.

\section{Methodology and Architecture}

\subsection{System Architecture}
The VeriByte toolchain is organized into modular phases forming an end-to-end compiler and verification boundary, depicted in Fig.~\ref{fig:arch}.

\begin{figure}[htbp]
\centering
\fbox{\parbox{0.95\columnwidth}{\centering
\small
\textbf{Source Program} (\texttt{.ml}) $\longrightarrow$ \textbf{Lexer \& Parser} (AST) $\longrightarrow$ \textbf{Semantic Analyzer}\\
$\downarrow$\\
\textbf{Bytecode Generator} $\longrightarrow$ [\textbf{Bytecode Optimizer}] $\longrightarrow$ \textbf{Serialized Binary} (\texttt{.bc})\\
$\downarrow$\\
\fbox{\textbf{STATIC VERIFIER} (Decoder $\to$ CFG $\to$ Worklist Abstract Interpreter)}\\
$\swarrow$ \quad \quad \quad \quad $\searrow$\\
\textit{Fail}: \textbf{Reject with Diagnostic} \quad \quad \textit{Pass}: \textbf{Admit to Bytecode VM}
}}
\caption{VeriByte End-to-End Pipeline Architecture.}
\label{fig:arch}
\end{figure}

The toolchain operates across two primary pipelines:
\begin{enumerate}
    \item \textbf{Source Compilation Pipeline:} Source files (\texttt{.ml}) pass through lexical analysis, recursive-descent parsing with precedence climbing, two-pass semantic type analysis, stack bytecode generation, and optional optimization before entering the verifier and VM.
    \item \textbf{Bytecode Admission Pipeline:} Externally supplied or precompiled bytecode files (\texttt{.bc}) enter directly at the serializer loader. The loader decodes header tables and immediately invokes the static verifier. Only bytecode certified as safe is permitted to initialize the virtual machine execution frames.
\end{enumerate}

\subsection{Source Language: MiniLang}
MiniLang is a strongly typed imperative language supporting 32-bit signed integers and booleans. Its syntax includes variable declarations (\texttt{let}), lexical block scopes, conditional branches (\texttt{if} / \texttt{else if} / \texttt{else}), iterative loops (\texttt{while}), first-class function declarations with explicit signatures, return statements, and console printing (\texttt{print}).

The semantic analyzer enforces:
\begin{itemize}
    \item Lexical scope isolation and duplicate declaration rejection.
    \item Strict type-checking across binary expressions ($+, -, *, /, \%, ==, !=, <, >, <=, >=, \&\&, ||$).
    \item Mandatory return paths on all control-flow trajectories for non-void functions.
    \item Bounds enforcement restricting integer literals strictly to $[-2^{31}, 2^{31}-1]$.
\end{itemize}

\subsection{Bytecode ISA and Binary File Format}
VeriByte targets a compact 27-opcode stack ISA, categorized in Table~\ref{tab:isa}. Multi-byte immediate operands are serialized in little-endian format.

\begin{table}[htbp]
\caption{VeriByte Bytecode Instruction Set Architecture (ISA)}
\label{tab:isa}
\centering
\begin{tabular}{llp{4.2cm}}
\toprule
\textbf{Category} & \textbf{Opcodes} & \textbf{Operand / Stack Transition} \\
\midrule
Stack Manipulation & \texttt{PUSH\_INT} $k$ & $\epsilon \to [\text{int}]$ (4-byte immediate) \\
                   & \texttt{PUSH\_BOOL} $b$ & $\epsilon \to [\text{bool}]$ (1-byte immediate) \\
                   & \texttt{POP} & $[T] \to \epsilon$ \\
                   & \texttt{DUP} & $[T] \to [T, T]$ \\
\midrule
Arithmetic (i32)   & \texttt{ADD, SUB, MUL} & $[\text{int}, \text{int}] \to [\text{int}]$ \\
                   & \texttt{DIV, MOD} & $[\text{int}, \text{int}] \to [\text{int}]$ (traps on 0) \\
                   & \texttt{NEG} & $[\text{int}] \to [\text{int}]$ \\
\midrule
Logic \& Compare   & \texttt{AND, OR} & $[\text{bool}, \text{bool}] \to [\text{bool}]$ \\
                   & \texttt{NOT} & $[\text{bool}] \to [\text{bool}]$ \\
                   & \texttt{EQ, NE} & $[T, T] \to [\text{bool}]$ (types must match) \\
                   & \texttt{LT, GT, LE, GE} & $[\text{int}, \text{int}] \to [\text{bool}]$ \\
\midrule
Local Variables    & \texttt{LOAD} $i$ & $\epsilon \to [T_i]$ (2-byte index) \\
                   & \texttt{STORE} $i$ & $[T_i] \to \epsilon$ (2-byte index) \\
\midrule
Control Flow       & \texttt{JMP} $addr$ & $\epsilon \to \epsilon$ (2-byte target offset) \\
                   & \texttt{JMP\_IF\_FALSE} $addr$ & $[\text{bool}] \to \epsilon$ (2-byte target offset) \\
\midrule
Functions \& I/O   & \texttt{CALL} $f$ & $[T_1, \dots, T_k] \to [T_{ret}]$ (2-byte fn) \\
                   & \texttt{RET} & $[T_{ret}] \to \text{caller stack}$ \\
                   & \texttt{PRINT} & $[T] \to \epsilon$ \\
                   & \texttt{HALT} & $[T] \to \text{exit}$ \\
\bottomrule
\end{tabular}
\end{table}

The serialized binary format (\texttt{.bc}) consists of an immutable header:
\begin{enumerate}
    \item \textbf{Magic Number \& Version:} 4-byte ASCII \texttt{SBVM} followed by a 2-byte little-endian version code (\texttt{0x0001}).
    \item \textbf{Function Table Header:} 2-byte function count and 2-byte entry function index.
    \item \textbf{Function Descriptors:} Per function: 1-byte name length, ASCII name, 1-byte parameter count, parameter type array, 1-byte return type, 2-byte local slot count, local type array, 2-byte declared maximum stack depth (\texttt{max\_stack}), 4-byte code length, and raw instruction bytes.
\end{enumerate}

\subsection{Static Bytecode Verifier (Core Contribution)}
The static verifier is constructed as an abstract interpreter operating over a monotone dataflow framework. For each function $F$, the verifier establishes:
\begin{equation}
\sigma \in \Sigma = \text{StackType}^* \times \text{LocalType}^M
\end{equation}
where $\text{StackType} \in \{\text{int}, \text{bool}\}$ and $\text{LocalType} \in \{\text{int}, \text{bool}, \text{UNSET}, \text{CONFLICT}\}$.

The verification algorithm proceeds in five sequential phases:

\subsubsection{Module Header Validation}
The module header is verified to guarantee structural consistency:
\begin{itemize}
    \item The module contains at least one function ($|F| \ge 1$).
    \item The entry function index references a valid function named \texttt{main} taking zero parameters and returning \texttt{int}.
    \item For every function, $\text{len}(\text{local\_types}) = \text{num\_locals}$, $\text{num\_locals} \ge \text{num\_params}$, and the prefix of $\text{local\_types}$ precisely matches $\text{param\_types}$.
\end{itemize}

\subsubsection{Linear Instruction Decoding}
The function byte stream is decoded sequentially. The verifier checks that every opcode is valid, operands do not truncate beyond code boundaries, and \texttt{PUSH\_BOOL} immediates are strictly \texttt{0x00} or \texttt{0x01}. The offset of every instruction boundary is recorded into an interval set $\mathcal{I}_{starts}$.

\subsubsection{Control-Flow Graph \& Jump Boundary Verification}
For every control-flow transfer instruction (\texttt{JMP}, \texttt{JMP\_IF\_FALSE}), the target address $A_{target}$ is validated against $\mathcal{I}_{starts}$:
\begin{equation}
A_{target} \in \mathcal{I}_{starts} \quad \land \quad 0 \le A_{target} < \text{len}(\text{code})
\end{equation}
Jumps landing mid-instruction or exceeding function code bounds are rejected immediately with \texttt{BAD\_JUMP}.

\subsubsection{Worklist Abstract Interpretation}
Verification initializes a double-ended worklist queue $\mathcal{W}$ with entry offset 0. The entry abstract state $\sigma_0$ is defined as:
\begin{equation}
\sigma_0 = \langle \text{stack} = [], \text{locals} = [T_1, \dots, T_k, \text{UNSET}, \dots, \text{UNSET}] \rangle
\end{equation}
where parameters $T_1 \dots T_k$ are typed, and all local variables allocated beyond parameters are marked \texttt{UNSET}.

For each instruction popped from $\mathcal{W}$, the verifier evaluates preconditions, computes the transfer function $\tau(op, \sigma) \to \sigma'$, and merges $\sigma'$ into the states of its control-flow successors.

\subsubsection{Join-Point Merging and Definite Assignment}
When multiple control-flow paths merge at offset $j$, the incoming state $\sigma_{in}$ is merged with the existing state $\sigma_{j}$ according to the confluence operator $\sqcap$:
\begin{align}
\text{depth}(\sigma_{in}.\text{stack}) &= \text{depth}(\sigma_j.\text{stack}) \quad (\text{else } \texttt{MERGE\_MISMATCH}) \\
\sigma_j.\text{stack}[k] &= \sigma_{in}.\text{stack}[k] \quad (\text{else } \texttt{MERGE\_MISMATCH}) \\
\sigma_j.\text{locals}[i] &= \begin{cases}
T & \text{if } \sigma_{in}.\text{locals}[i] = \sigma_j.\text{locals}[i] = T \\
\text{UNSET} & \text{if either is } \text{UNSET} \\
\text{CONFLICT} & \text{otherwise}
\end{cases}
\end{align}
If local $i$ is not initialized along all convergent execution trajectories, it resolves to \texttt{UNSET}. Any subsequent \texttt{LOAD} $i$ instruction encountering an \texttt{UNSET} state is rejected with \texttt{LOCAL\_UNSET}.

\subsubsection{Comprehensive Invariant Taxonomy}
The verifier strictly enforces the thirteen invariant rules specified in Table~\ref{tab:rules}.

\begin{table}[htbp]
\caption{The Thirteen VeriByte Static Verification Rules}
\label{tab:rules}
\centering
\begin{tabular}{lp{6.2cm}}
\toprule
\textbf{Rule Name} & \textbf{Violation Condition \& Security Rationale} \\
\midrule
\texttt{BAD\_OPCODE} & Undefined byte value or invalid boolean immediate. \\
\texttt{TRUNCATED} & Byte stream terminates before reading immediate operand. \\
\texttt{BAD\_JUMP} & Jump target is out of bounds or lands mid-instruction. \\
\texttt{STACK\_UNDERFLOW} & Popping when stack depth is insufficient for opcode. \\
\texttt{STACK\_OVERFLOW} & Pushing when stack depth equals declared \texttt{max\_stack}. \\
\texttt{STACK\_TYPE} & Operand type does not match required instruction type. \\
\texttt{MERGE\_MISMATCH} & Converging branches have different stack depths or types. \\
\texttt{LOCAL\_RANGE} & Accessing local index $\ge \text{num\_locals}$. \\
\texttt{LOCAL\_TYPE} & Storing an operand incompatible with declared local type. \\
\texttt{LOCAL\_UNSET} & Reading a local variable before definite assignment. \\
\texttt{BAD\_CALL} & Callee index invalid or argument types mismatch signature. \\
\texttt{BAD\_RETURN} & Stack depth $\ne 1$ or top value mismatches function return type. \\
\texttt{FALLTHROUGH} & Instruction flow runs off end of code without \texttt{RET}/\texttt{HALT}. \\
\bottomrule
\end{tabular}
\end{table}

\subsection{Virtual Machine Interpreter \& Checked Soundness Oracle}
The virtual machine executes verified bytecode through an activation-frame stack. Each call frame allocates an isolated local variable array and an operand stack. Arithmetic instructions evaluate over 32-bit two's complement representation:
\begin{equation}
\text{to\_i32}(v) = (v + 2^{31} \pmod{2^{32}}) - 2^{31}
\end{equation}
Division or modulo by zero raises a clean, recoverable \texttt{VMTrap("division by zero")} without crashing the host process. Call depth recursion is bounded at 10,000 frames to prevent host stack exhaustion.

To validate soundness during fuzzing, the VM incorporates a \textbf{Checked Mode} (\texttt{checked=True}). In checked mode, the VM validates hardware-level safety invariants at every cycle:
\begin{enumerate}
    \item Arithmetic/comparison operands must satisfy $\text{type}(v) \equiv \text{int}$.
    \item Logical and conditional branch operands must satisfy $\text{type}(v) \equiv \text{bool}$.
    \item \texttt{EQ}/\texttt{NE} operands must have identical concrete types.
    \item \texttt{LOAD} must not yield \texttt{None} (confirming definite assignment).
    \item \texttt{STORE} must match declared slot types.
    \item Stack depth after any push must not exceed declared \texttt{max\_stack}.
    \item \texttt{RET} must leave exactly one value matching function return type.
    \item Program counter must remain within valid instruction start boundaries.
\end{enumerate}
Any invariant violation in checked mode raises \texttt{VMInvariantError}, flagging a soundness defect in the static verifier.

\subsection{Bytecode Optimization Pipeline}
The bytecode optimizer lifts raw bytecode into an Intermediate Representation (IR) consisting of explicit basic blocks. Optimization passes include:
\begin{itemize}
    \item \textbf{Constant Folding:} Algebraic reduction of constant arithmetic and comparison operations (e.g., $10 + 20 \to 30$).
    \item \textbf{Dead Instruction Elimination:} Removal of redundant \texttt{PUSH} followed immediately by \texttt{POP}, and double logical negation (\texttt{NOT NOT}).
    \item \textbf{Jump Threading:} Elimination of jumps pointing to the immediate subsequent instruction offset.
    \item \textbf{Unreachable Code Elimination:} Dead block pruning following unconditional returns or jumps.
\end{itemize}
Crucially, VeriByte enforces verified optimization: the pipeline verifies the input bytecode, executes optimization passes, recalculates jump targets and stack limits, and re-verifies the optimized bytecode prior to execution.

\section{Results and Discussion}

The experimental evaluation of VeriByte investigates five empirical research questions:
\begin{itemize}
    \item \textbf{RQ1 (Precision):} Does the verifier accept all valid compiled programs with zero false rejections?
    \item \textbf{RQ2 (Attack Rejection):} Does the verifier accurately detect and isolate each of the 13 safety violations?
    \item \textbf{RQ3 (Soundness):} Across random differential fuzzing, do accepted mutants execute without invariant violations?
    \item \textbf{RQ4 (Loader Hardening):} How robust is binary deserialization against raw bit flips?
    \item \textbf{RQ5 (Performance):} What is the computational overhead and scaling factor of load-time verification?
\end{itemize}

\subsection{Unit Test Suite and Precision (RQ1)}
The regression test suite consists of 195 automated tests executing via the Python standard library test runner. The test suite covers:
\begin{itemize}
    \item Front-end lexing and parsing: 22 tests.
    \item Semantic type analysis and bounds checking: 56 tests.
    \item Code generation, serialization, and disassembly: 13 tests.
    \item Virtual machine execution and wraparound arithmetic: 18 tests.
    \item Static verifier rule isolation: 38 tests.
    \item Disk-based malicious binary suite: 1 test (26 subtests).
    \item Bytecode optimizer equivalence: 13 tests.
    \item Malformed binary and literal hardening: 7 tests.
    \item Checked-mode soundness invariants and oracle meta-tests: 27 tests.
\end{itemize}
All 195 tests pass with zero failures or errors. Across all canonical example programs (\texttt{factorial.ml}, \texttt{fibonacci.ml}, \texttt{gcd.ml}, \texttt{logic.ml}, \texttt{prime.ml}), verification achieved a \textbf{0\% false-rejection rate} both before and after bytecode optimization.

\subsection{Handcrafted Exploit Corpus Validation (RQ2)}
To evaluate rule precision, we synthesized 26 malicious bytecode files (two distinct files per safety rule) stored permanently on disk in \texttt{tests/malicious/}. Table~\ref{tab:malicious_results} details the evaluation results.

\begin{table}[htbp]
\caption{Malicious Exploit Corpus Evaluation}
\label{tab:malicious_results}
\centering
\begin{tabular}{lllcc}
\toprule
\textbf{Target Exploit File} & \textbf{Target Rule} & \textbf{Observed Rule} & \textbf{Unsafe VM} & \textbf{Verified VM} \\
\midrule
\texttt{bad\_opcode\_1, 2.bc} & \texttt{BAD\_OPCODE} & \texttt{BAD\_OPCODE} & Crash & Rejected \\
\texttt{truncated\_1, 2.bc} & \texttt{TRUNCATED} & \texttt{TRUNCATED} & Crash & Rejected \\
\texttt{bad\_jump\_1, 2.bc} & \texttt{BAD\_JUMP} & \texttt{BAD\_JUMP} & Crash & Rejected \\
\texttt{stack\_underflow\_1, 2.bc} & \texttt{STACK\_UNDERFLOW} & \texttt{STACK\_UNDERFLOW} & Crash & Rejected \\
\texttt{stack\_overflow\_1, 2.bc} & \texttt{STACK\_OVERFLOW} & \texttt{STACK\_OVERFLOW} & Silent & Rejected \\
\texttt{stack\_type\_1, 2.bc} & \texttt{STACK\_TYPE} & \texttt{STACK\_TYPE} & Confused & Rejected \\
\texttt{merge\_mismatch\_1, 2.bc} & \texttt{MERGE\_MISMATCH} & \texttt{MERGE\_MISMATCH} & Crash & Rejected \\
\texttt{local\_range\_1, 2.bc} & \texttt{LOCAL\_RANGE} & \texttt{LOCAL\_RANGE} & Crash & Rejected \\
\texttt{local\_type\_1, 2.bc} & \texttt{LOCAL\_TYPE} & \texttt{LOCAL\_TYPE} & Confused & Rejected \\
\texttt{local\_unset\_1, 2.bc} & \texttt{LOCAL\_UNSET} & \texttt{LOCAL\_UNSET} & Crash & Rejected \\
\texttt{bad\_call\_1, 2.bc} & \texttt{BAD\_CALL} & \texttt{BAD\_CALL} & Crash & Rejected \\
\texttt{bad\_return\_1, 2.bc} & \texttt{BAD\_RETURN} & \texttt{BAD\_RETURN} & Crash & Rejected \\
\texttt{fallthrough\_1, 2.bc} & \texttt{FALLTHROUGH} & \texttt{FALLTHROUGH} & Crash & Rejected \\
\bottomrule
\end{tabular}
\end{table}

When executed under unverified mode (\texttt{--unsafe}), 22 of the 26 binaries crashed catastrophically with unhandled host exceptions (\texttt{IndexError: pop from empty list}, \texttt{IndexError: list index out of range}), while the remaining 4 exhibited severe type confusion or silent memory overconsumption. Under verified mode, \textbf{100\% of the malicious binaries were cleanly intercepted and rejected} with precise diagnostic error messages detailing the function name, instruction offset, mnemonic, violated rule, and stack state.

\subsection{Differential Mutation Fuzzing (RQ3)}
To test verifier soundness, we engineered an automated mutation fuzzer applying 11 mutation operators spanning instruction corruption (bit flips, byte substitution, opcode swaps, truncation, byte insertion/deletion) and metadata corruption (stack limits, local counts, entry function index, parameter/return/local types).

We executed a campaign of 15,000 mutants across the 5 canonical seed programs over three independent random seeds (101, 202, 303), testing 1,000 mutations per program per seed with bytecode optimization (\texttt{--opt}) and active VM checked mode. Table~\ref{tab:fuzz_summary} summarizes the results.

\begin{table}[htbp]
\caption{15,000-Mutant Differential Fuzzing Campaign Results}
\label{tab:fuzz_summary}
\centering
\begin{tabular}{lrrrr}
\toprule
\textbf{Metric} & \textbf{Seed 101} & \textbf{Seed 202} & \textbf{Seed 303} & \textbf{Total / Overall} \\
\midrule
Total Mutants Tested & 5,000 & 5,000 & 5,000 & \textbf{15,000} \\
Accepted by Verifier & 945 (18.9\%) & 882 (17.6\%) & 942 (18.8\%) & \textbf{2,769 (18.5\%)} \\
Rejected by Verifier & 4,055 (81.1\%) & 4,118 (82.4\%) & 4,058 (81.2\%) & \textbf{12,231 (81.5\%)} \\
Unsafe Mode Crashes  & 3,154 & 3,144 & 3,123 & \textbf{9,421 (77.0\%)} \\
\textbf{Soundness Violations} & \textbf{0} & \textbf{0} & \textbf{0} & \textbf{0 (Zero)} \\
\bottomrule
\end{tabular}
\end{table}

The rejection distribution across the 12,231 rejected mutants exercised all 13 safety rules:
\begin{itemize}
    \item \texttt{BAD\_OPCODE}: 5,402 (44.2\%)
    \item \texttt{LOCAL\_RANGE}: 1,449 (11.8\%)
    \item \texttt{BAD\_CALL}: 1,200 (9.8\%)
    \item \texttt{TRUNCATED}: 1,144 (9.4\%)
    \item \texttt{BAD\_JUMP}: 667 (5.5\%)
    \item \texttt{FALLTHROUGH}: 632 (5.2\%)
    \item \texttt{STACK\_OVERFLOW}: 481 (3.9\%)
    \item \texttt{STACK\_UNDERFLOW}: 397 (3.2\%)
    \item \texttt{BAD\_RETURN}: 340 (2.8\%)
    \item \texttt{LOCAL\_TYPE}: 240 (2.0\%)
    \item \texttt{STACK\_TYPE}: 141 (1.2\%)
    \item \texttt{LOCAL\_UNSET}: 112 (0.9\%)
    \item \texttt{MERGE\_MISMATCH}: 26 (0.2\%)
\end{itemize}

Crucially, across all 2,769 accepted mutants executed in checked mode, \textbf{zero soundness violations occurred}. 2,528 programs ran to clean completion, 43 terminated safely via expected division-by-zero traps (\texttt{VMTrap}), and 198 safely terminated via step-limit timeouts. When rejected mutants were executed in unverified mode, \textbf{9,421 crashed with unhandled host exceptions}, demonstrating that the verifier successfully prevented thousands of potential runtime faults.

\subsection{Raw-Byte Serializer Hardening Fuzzing (RQ4)}
To evaluate whether corrupted binary files could trigger unhandled exceptions in the binary loader, we executed the fuzzer in file-level mode (\texttt{--file-mode}), mutating raw serialized \texttt{.bc} bytes directly across 10,000 mutants (Seeds 101 and 202). Table~\ref{tab:file_fuzz} outlines the findings.

\begin{table}[htbp]
\caption{10,000-Mutant Raw Binary Deserializer Fuzzing}
\label{tab:file_fuzz}
\centering
\begin{tabular}{lrrr}
\toprule
\textbf{Outcome} & \textbf{Seed 101} & \textbf{Seed 202} & \textbf{Combined} \\
\midrule
Total Mutated Files & 5,000 & 5,000 & 10,000 \\
Loader Rejections (\texttt{FORMAT\_ERROR}) & 2,787 (55.7\%) & 2,791 (55.8\%) & 5,578 (55.8\%) \\
Verifier Rejections (\texttt{VerifyError}) & 1,592 (31.8\%) & 1,570 (31.4\%) & 3,162 (31.6\%) \\
Accepted Files & 621 (12.4\%) & 639 (12.8\%) & 1,260 (12.6\%) \\
\textbf{Soundness Violations} & \textbf{0} & \textbf{0} & \textbf{0 (Zero)} \\
\bottomrule
\end{tabular}
\end{table}

Across 10,000 raw corrupted files, exactly 5,578 were intercepted by the hardened binary loader with clean \texttt{BytecodeFormatError} exceptions, and 3,162 passed the loader but were rejected by the static verifier. Zero unhandled Python exceptions escaped, demonstrating complete defensive encapsulation of the binary ingestion boundary.

\subsection{Validating the Soundness Oracle (Meta-Test)}
A potential threat to validity in fuzzing campaigns is a vacuous oracle—an oracle that never reports a failure because its internal invariant assertions are disabled. To validate our oracle's sensitivity, we implemented an automated meta-test in \texttt{tests/test_oracle.py}.

The test temporarily disables the verifier's operand-type check specifically for \texttt{Op.ADD} (\texttt{CHECK\_STACK\_TYPE\_ADD = False}) and synthesizes a program pushing a boolean and an integer before executing \texttt{ADD}. The weakened verifier improperly accepts the binary. Upon admission to the checked VM, the runtime oracle immediately detects the non-integer operand and raises \texttt{VMInvariantError}, recording a soundness violation. This meta-test confirms that the checked VM oracle reliably identifies real verifier omissions.

\subsection{Performance Benchmarks and Scaling (RQ5)}
Verification performance was profiled across synthetic programs of increasing size (up to 1,000 statements / 12,012 bytecode bytes). Fig.~\ref{fig:scaling} illustrates verification latency as a function of bytecode size.

\begin{figure}[htbp]
\centering
\fbox{\parbox{0.88\columnwidth}{\small
\textbf{Verification Latency vs. Program Size:}\\
$\bullet$ 252 Bytes (20 stmts) : $\approx 1.4\text{ ms}$ ($176\text{ KB/s}$)\\
$\bullet$ 612 Bytes (50 stmts) : $\approx 3.0\text{ ms}$ ($198\text{ KB/s}$)\\
$\bullet$ 1,212 Bytes (100 stmts): $\approx 7.1\text{ ms}$ ($167\text{ KB/s}$)\\
$\bullet$ 3,012 Bytes (250 stmts): $\approx 15.9\text{ ms}$ ($185\text{ KB/s}$)\\
$\bullet$ 6,012 Bytes (500 stmts): $\approx 31.2\text{ ms}$ ($188\text{ KB/s}$)\\
$\bullet$ 12,012 Bytes (1,000 stmts): $\approx 66.3\text{ ms}$ ($177\text{ KB/s}$)\\
\vspace{1mm}
\textbf{Empirical Complexity: Strictly Linear $\mathcal{O}(N)$ Scaling}
}}
\caption{Verification Latency and Throughput Scaling.}
\label{fig:scaling}
\end{figure}

Throughput remained steady at approximately $180\text{ KB/s}$ across all workloads. When evaluated against total application execution time across realistic benchmarks (\texttt{factorial}, \texttt{fibonacci}, \texttt{gcd}, \texttt{prime}), static verification accounted for \textbf{less than 0.5\% of runtime execution latency} (and under 0.1\% for compute-heavy recursive workloads). Bytecode optimization yielded up to \textbf{28.9\% bytecode size reduction} and up to \textbf{1.28x execution speedup} on constant-heavy routines.

\section{Conclusions and Future Work}
This paper presented VeriByte, a secure bytecode compilation, verification, and execution toolchain. By elevating the compiler and virtual machine loader into a formal trust boundary, VeriByte guarantees that untrusted bytecode cannot induce stack underflows, type confusion, invalid control-flow transfers, or uninitialized memory reads.

Through abstract interpretation over a monotone type lattice, VeriByte enforces thirteen structural and behavioral safety rules. Across 195 unit tests, 26 handcrafted exploit binaries, 15,000 structural mutants, and 10,000 raw binary byte mutations, the static verifier achieved zero false rejections on valid programs, intercepted over 9,400 unhandled execution crashes, and recorded zero soundness violations under active runtime oracle checking. With strictly linear $\mathcal{O}(N)$ scaling and an execution overhead below 0.5\%, VeriByte demonstrates that static verification provides robust, cost-effective security for intermediate code execution.

\subsection{Future Work}
Future research directions include:
\begin{enumerate}
    \item \textbf{Array and Heap Allocation Bounds:} Extending the type system and verifier to support fixed-size arrays with static index analysis and provably safe heap references.
    \item \textbf{Stack Map Tables:} Transitioning from full worklist abstract interpretation to single-pass linear verification via compiler-generated stack map tables, mirroring modern Java 6+ class file verification.
    \item \textbf{Short-Circuit Logic Evaluation:} Extending the code generator and AST lowering to support short-circuit boolean semantics for \texttt{\&\&} and \texttt{||}.
    \item \textbf{Formal Machine-Checked Soundness Proofs:} Formalizing the operational semantics and progress-preservation theorems in interactive theorem provers such as Coq or Lean.
    \item \textbf{Just-In-Time (JIT) Compilation:} Integrating a native JIT compiler (targeting x86-64 or ARM64) that directly compiles verified bytecode with eliminated runtime safety checks.
\end{enumerate}

\begin{thebibliography}{00}
\bibitem{lindholm2014java} T. Lindholm, F. Yellin, B. Bracha, and A. Buckley, \textit{The Java Virtual Machine Specification}, Java SE 8 ed. Boston, MA, USA: Addison-Wesley, 2014.
\bibitem{meijer2002technical} E. Meijer and J. Gough, ``Technical overview of the Common Language Runtime,'' \textit{Microsoft Research}, Tech. Rep., 2002.
\bibitem{haas2017bringing} A. Haas, A. Rossberg, D. L. Schuff, B. L. Titzer, M. Holman, D. Gohman, L. Wagner, A. Zakai, and J. F. Bastien, ``Bringing the web up to speed with WebAssembly,'' in \textit{Proc. 38th ACM SIGPLAN Conf. Program. Lang. Des. Implement. (PLDI)}, 2017, pp. 185--200.
\bibitem{ebpf2023linux} Linux Kernel Documentation, ``The eBPF verifier,'' kernel.org, 2023. [Online]. Available: \url{https://docs.kernel.org/bpf/verifier.html}
\bibitem{gershuni2019simple} E. Gershuni, N. Amit, A. Gurfinkel, N. Narodytska, J. A. Navas, N. Rinetzky, L. Ryzhyk, and M. Sagiv, ``Simple and precise static analysis of untrusted Linux kernel extensions,'' in \textit{Proc. 40th ACM SIGPLAN Conf. Program. Lang. Des. Implement. (PLDI)}, 2019, pp. 1069--1084.
\bibitem{kildall1973unified} G. A. Kildall, ``A unified approach to global program optimization,'' in \textit{Proc. 1st Annu. ACM SIGACT-SIGPLAN Symp. Princ. Program. Lang. (POPL)}, 1973, pp. 194--206.
\bibitem{cousot1977abstract} P. Cousot and R. Cousot, ``Abstract interpretation: A unified lattice model for static analysis of programs by construction or approximation of fixpoints,'' in \textit{Proc. 4th ACM SIGACT-SIGPLAN Symp. Princ. Program. Lang. (POPL)}, 1977, pp. 238--252.
\bibitem{leroy2003java} X. Leroy, ``Java bytecode verification: An overview,'' \textit{J. Autom. Reason.}, vol. 30, no. 3--4, pp. 235--260, 2003.
\bibitem{aho2006compilers} A. V. Aho, M. S. Lam, R. Sethi, and J. D. Ullman, \textit{Compilers: Principles, Techniques, and Tools}, 2nd ed. Boston, MA, USA: Addison-Wesley, 2006.
\bibitem{nystrom2021crafting} R. Nystrom, \textit{Crafting Interpreters}. Seattle, WA, USA: Genever Benning, 2021.
\end{thebibliography}

\end{document}
